%% DILARANG EDIT BAGIAN INI
\clearpage
\phantomsection
\addcontentsline{toc}{chapter}{KATA PENGANTAR}
\begin{center}
    \textbf{\large KATA PENGANTAR}\\[3em]
\end{center}
%% DILARANG EDIT BAGIAN INI

Segala puji beserta syukur senantiasa penulis panjatkan kepada Allah SWT yang telah memberikan rahmat dan petunjuk-Nya kepada penulis sehingga dapat menyelesaikan proposal skripsi yang berjudul “Rancang Bangun Aplikasi Point of Sale dengan Fitur Prediksi Penjualan Harian Menggunakan Metode Long Short-Term Memory (LSTM)”. Shalawat beserta salam kepada baginda Rasulullah Muhammad SAW, keluarga, dan para sahabat yang telah berjasa dalam memperjuangkan islam dan ilmu pengetahuan. Proposal skripsi ini disusun sebagai salah satu syarat untuk menyelesaikan pendidikan Sarjana Terapan Program Studi Teknik Informatika, Jurusan Teknologi Informasi dan Komputer, Politeknik Negeri Lhokseumawe. 

Dalam Penyusunan skripsi ini penulis banyak mendapatkan dukungan, perhatian, dorongan, bimbingan dan berbagai bantuan dari banyak pihak terutama para dosen pembimbing. Dengan demikian, segala hormat dan kerendahan hati penulis ini menyampaikan rasa terima kasih yang tidak terhingga terutama kepada:
\begin{enumerate}
    \item Allah SWT, atas segala rahmat, karunia, dan hidayah-Nya yang telah diberikan sehingga penulis dapat menyelesaikan skripsi ini.
    \item Kedua orang tua tercinta serta keluarga yang senantiasa memberikan doa, dukungan, kasih sayang, dan pengorbanan tanpa henti kepada penulis, sehingga penulis dapat menyelesaikan pendidikan dan penyusunan skripsi ini dengan lancar.
    \item Bapak Ir. Rizal Syahyadi, S.T., M.Eng.Sc., IPM., ASEAN Eng., APEC Eng. selaku Direktur Politeknik Negeri Lhokseumawe.
    \item Bapak Salahuddin, S.T., M.Cs selaku Ketua Jurusan Teknologi Informasi dan Komputer, Bapak Fachri Yanuar Rudi F, M.T selaku Sekretaris Jurusan Teknologi Informasi dan Komputer, dan Bapak M. Khadafi, S.T., M.T selaku Ketua Program Studi Teknik Informatika atas segala bimbingan dan perhatian selama masa perkuliahan.
    \item Bapak Zulfan Khairil Simbolon, S.T., M.Eng. selaku Pembimbing utama dan Bapak Azhar, S.T., M.T. selaku pembimbing pendamping yang telah membantu, membimbing, dan mengarahkan penulis dalam proses penyelesaian proposal skripsi ini.
    \item Para dosen dan tenaga pengajar Program Studi Teknik Informatika yang telah memberikan ilmu, wawasan, serta pengalaman berharga kepada penulis selama masa perkuliahan, sehingga menjadi bekal yang sangat berarti dalam proses penyelesaian studi.
    \item Seluruh teman-teman, khususnya teman-teman terdekat, yang selalu menemani dan memberikan dukungan selama masa perkuliahan hingga penyusunan skripsi ini. Terima kasih atas kebersamaan, perhatian, dan semangat yang tidak pernah berhenti diberikan kepada penulis. Meskipun tidak dapat disebutkan satu per satu, penulis sangat menghargai setiap bantuan dan dukungan yang telah diberikan.
    \item Kepada diri sendiri, penulis mengucapkan terima kasih atas segala usaha, kesabaran, dan kekuatan dalam melewati setiap proses yang tidak mudah. Di tengah keraguan dan kelelahan, penulis tetap berusaha bertahan dan melangkah hingga tahap akhir. Skripsi ini menjadi bukti dari perjalanan panjang yang dilalui serta hasil dari usaha dan doa yang terus diupayakan.
\end{enumerate}

Penulis menyadari bahwa skripsi ini masih jauh dari sempurna dan memiliki berbagai keterbatasan. Oleh karena itu, penulis sangat mengharapkan kritik dan saran yang bersifat membangun demi perbaikan di masa yang akan datang. Semoga segala bantuan dan kebaikan yang telah diberikan mendapatkan balasan yang terbaik.

%% DILARANG EDIT BAGIAN INI
\begin{flushright}
    Lhokseumawe, \tglpernyataan\\[1.25cm]
    \penulis \\
    NIM. \nim
\end{flushright}
%% DILARANG EDIT BAGIAN INI
